\section{Conclusiones}
\label{sec:conclusiones}

Para la fundación estudiada, la colocación continua a $60\,\mathrm{m^3/h}$ terminó en $10{,}61\,\mathrm h$ y cumplió los criterios térmicos en los $24$ escenarios. Frente al plan de $55\,\mathrm{m^3/h}$ ahorró $57{,}89\,\mathrm{min}$, equivalentes al $8{,}33\,\text{\%}$ de la duración de la alternativa más lenta. Al sostener dos frentes de $30\,\mathrm{m^3/h}$ sin esperas, la propuesta alcanzó la menor duración compatible con el volumen y la capacidad disponible.

La probabilidad conjunta calculada fue $0\,\text{\%}$ tanto a los $14$ como a los $28$ días. En la malla fina, la mayor temperatura fue $47{,}03\,^\circ\mathrm C$ y el mayor diferencial fue $15{,}72\,\mathrm K$, con un margen de $4{,}28\,\mathrm K$ respecto del criterio más cercano. Estos resultados respaldan la hipótesis de trabajo para el catálogo estudiado y la programación que mantiene la última manta hasta los $21$ días de edad local y la sombra durante el seguimiento.

La principal condición para trasladar esta decisión a obra es comprobar que se mantengan las prestaciones simuladas, especialmente el ingreso del hormigón a $12\,^\circ\mathrm C$, la extracción de los nueve circuitos y la atenuación solar de la cubierta. Las mediciones de núcleo, contorno y ambiente permitirían contrastar la predicción y fundamentar el retiro de cada protección. Quedan pendientes la validación térmica de la fundación y la evaluación de la retirada definitiva de la sombra y la refrigeración.
